% Extensión estimada: unas 3.5 páginas, 10 párrafos y tres subsecciones (ficha 3.5).
% Fuentes: ficha fichas/03-materiales-metodos/05-clasificacion-vlm.md; config.json de Phi-3.5-vision-instruct;
% src/models/mini_vlm.py, llava_med.py, dermlip.py; src/evaluation/vlm_scoring.py; scripts/train_vlm_classify.py,
% embeddings_classical.py, dermlip_split_eval.py; docs/resul/split_lesion_resul_20260926.md (§2, §6, §8).
% Revisión del autor del 18/10/2026: los modelos se nombran al final de la introducción; transición a DermLIP.
\section{Clasificación con modelos visión--lenguaje}
\label{sec:met_clasif_vlm}

La primera pregunta de investigación es si un \gls{vlm} clasifica las lesiones tan bien como los clasificadores de
la sección anterior. Para responderla se usan dos tipos de \glspl{vlm} (Sección~\ref{sec:marco_vlm}). Los generativos
responden con texto, así que hay que decidir cómo se convierte una respuesta escrita en una de las siete clases. Los
contrastivos, en cambio, no generan texto, sino una representación de la imagen, sobre la que se puede clasificar directamente.
En concreto, se evalúan dos generativos que no están especializados en dermatología, uno de dominio general y otro biomédico, y
un contrastivo que sí lo está. Los generativos son Mini-VLM y LLaVA-Med.

\subsection{Mini-VLM y LLaVA-Med}
\label{sec:met_vlm_generativos}

Mini-VLM es el modelo Phi-3.5-vision-instruct~\cite{abdin2024phi3}, un \gls{vlm} de propósito general. Combina un
codificador de imagen \gls{clip} ViT-L/14 a 336 píxeles con una proyección de tipo MLP y un \gls{llm} de solo
decodificador, Phi-3.5-mini, de 32 capas y 3\,072 dimensiones, para un total de unos 4\,150 millones de parámetros.
Se eligió por dos motivos: sus pesos son abiertos y su tamaño permite aplicarle \gls{lora} en una GPU de consumo.

El segundo generador, LLaVA-Med~\cite{li2023llavamed}, en su versión 1.5 con Mistral-7B~\cite{zeng2025mmskin}, es el \gls{vlm} biomédico con el que se contrasta
Mini-VLM. Usa el mismo tipo de codificador, \gls{clip} ViT-L/14, una proyección y el modelo de lenguaje
Mistral-7B, pero se entrenó con figuras y leyendas de la literatura biomédica. Además, con un modelo de lenguaje de unos 7 000
millones de parámetros, es el más pesado de los dos. De este modo, la comparación pone frente a frente un modelo
entrenado con datos biomédicos y otro de propósito general.

Ninguno de los dos devuelve siempre una etiqueta válida cuando se le pide un diagnóstico, así que hace falta fijar
cómo se obtiene una clase con cada uno.

\subsection{Cuatro estrategias para obtener una clase}
\label{sec:met_vlm_estrategias}

Las cuatro estrategias usan el mismo conjunto de siete categorías y se diferencian en cómo se obtiene la clase. La
Figura~\ref{fig:met_zsvslik} contrasta las dos primeras: la que lee la respuesta frente a la que puntúa las
clases. Los \emph{prompts} completos están en el Apéndice~\ref{ap:prompts}.

\begin{enumerate}[leftmargin=1.6em]
  \item \textbf{\emph{Zero-shot} por generación.} El modelo recibe la imagen y un prompt con las clases posibles, y la
  clase se extrae del texto que genera con decodificación voraz (Sección~\ref{sec:marco_llm}). Se acepta la coincidencia
  exacta con una etiqueta o la primera que aparece como palabra completa. Si no hay etiqueta válida, la respuesta cuenta
  como «desconocido».
  \item \textbf{\emph{Likelihood scoring}.} En lugar de leer la respuesta, se compara la probabilidad que el modelo asigna al
  nombre de cada clase dada la imagen. La puntuación se define en la Ecuación~\ref{eq:met_likelihood}.
  \item \textbf{Razonamiento paso a paso.} Es la estrategia de \gls{cot}: el prompt pide describir los rasgos de la
  lesión, listar diagnósticos diferenciales y concluir con una línea «Final diagnosis:» seguida de la clase. Si esa línea
  falta, se usa el parseo de la primera estrategia, y sin etiqueta válida la respuesta cuenta como «desconocido».
  Además, la respuesta se limita a 32 \emph{tokens} nuevos en la primera estrategia y a 200 en esta, que necesita espacio para el
  razonamiento.
  \item \textbf{Ajuste fino con \gls{lora}.} Se evalúa el \emph{checkpoint} de la etapa 2 de generación, ajustado con LoRA~\cite{hu2021lora}
  sobre \emph{SkinCaRe}~\cite{shen2025skincare}, por verosimilitud o por \gls{cot}.
  Este checkpoint no se ajustó sobre \emph{HAM10000}. El procedimiento en dos etapas se describe en la
  Sección~\ref{sec:met_generacion}.
\end{enumerate}

% Figura 3.6: zero-shot frente a likelihood scoring. Adaptada de fig:zsvslik de Docs/experimentacion/experimentacion.tex,
% redibujada con los estilos de Configurations/13-Figuras.sty, sin cifras.
\begin{figure}[H]
\centering
\begin{tikzpicture}[node distance=0.6cm and 0.55cm]
  % (a) zero-shot: se lee la clase de la respuesta generada
  \node[font=\small, anchor=south west] (la) at (-1.45,0.95) {(a) \emph{zero-shot}};
  \node[figentrada, text width=2.7cm] (ai) at (0,0) {imagen + \emph{prompt}};
  \node[figtarea, text width=2.7cm, right=of ai] (ag) {el VLM genera texto libre};
  \node[figsalida, text width=2.9cm, right=of ag] (ap) {extracción de la clase o «desconocido»};
  \draw[figflecha] (ai) -- (ag);
  \draw[figflecha] (ag) -- (ap);

  % (b) likelihood: se puntúa el nombre de cada clase
  \node[font=\small, anchor=south west] (lb) at (-1.45,-1.75) {(b) \emph{likelihood}};
  \node[figentrada, text width=2.7cm] (bi) at (0,-2.95) {imagen + \emph{prompt}};
  \node[figtarea, text width=2.7cm, right=of bi] (bs) {el VLM puntúa $P(\text{nombre}_c \mid x)$ para cada clase $c$};
  \node[figsalida, text width=2.9cm, right=of bs] (ba) {$\arg\max_c$, sin «desconocido»};
  \draw[figflecha] (bi) -- (bs);
  \draw[figflecha] (bs) -- (ba);
\end{tikzpicture}
\caption{Dos formas de clasificar con un VLM generativo: leer la clase de la respuesta (a) o elegir la clase de
mayor puntaje (b).}
\label{fig:met_zsvslik}
\end{figure}


En la estrategia de verosimilitud, para cada imagen $x$ y cada clase $c$ se toman los tres sinónimos del nombre de la
clase. Con ellos, se calcula la log-probabilidad de cada sinónimo como continuación del prompt dada la imagen, se
divide por su cantidad de tokens para no favorecer los nombres cortos y se queda con el máximo entre ellos. Luego, ese
valor se promedia sobre los cuatro prompts y la clase predicha es la de mayor puntaje:
\begin{equation}
    s_c = \frac{1}{|\mathcal{P}|} \sum_{p \in \mathcal{P}} \max_{s \in S_c} \frac{\log P(s \mid p, x)}{|s|},
    \qquad \hat{c} = \arg\max_{c} s_c,
    \label{eq:met_likelihood}
\end{equation}
donde $\mathcal{P}$ es el conjunto de los cuatro prompts, $S_c$ los sinónimos de la clase $c$, $P(s \mid p, x)$ la
probabilidad que el modelo asigna a $s$ como continuación de $p$ dada la imagen $x$, como en la
Ecuación~\ref{eq:marco_lm}, y $|s|$ la cantidad de tokens de $s$. Como la comparación se hace entre clases, nunca produce
«desconocido». La respuesta generada, en cambio, depende de la decodificación y puede no contener ninguna etiqueta.

Una vez definido el puntaje, para comprobar si depende de la imagen, el mismo cálculo se repite con cada imagen
reemplazada por una gris neutra de $224 \times 224$ píxeles, con el valor 128 en cada canal. Si el modelo sigue dando la misma distribución de
puntajes, el prompt y no la imagen guía la respuesta. Este control se describe aquí como procedimiento, mientras
que sus resultados se presentan en los capítulos de resultados.

\subsection{DermLIP y sus modos de uso}
\label{sec:met_dermlip}

Con Mini-VLM y LLaVA-Med, la clase sale siempre de un texto, sea leyéndolo o puntuándolo. DermLIP~\cite{yan2025derm1miccv}
evita ese paso: es un \gls{vlm} contrastivo al estilo de \gls{clip}, con un codificador de imagen ViT-B/16 y uno de
texto, preentrenado sobre los pares dermatológicos de \emph{Derm1M} (Sección~\ref{sec:met_datasets}). No genera texto,
sino un \emph{embedding} de 512 dimensiones por imagen. Se incluye porque permite distinguir dos explicaciones de un mal
resultado de los \glspl{vlm} generales: que les falte especialización en dermatología, o que generar texto sea una mala
forma de clasificar.

El \emph{embedding} de cada imagen, congelado y normalizado, se usa de cuatro maneras:
\begin{enumerate}[leftmargin=1.6em]
  \item \textbf{Zero-shot.} Se compara el \emph{embedding} de la imagen con el de cada clase por similitud coseno
  (Ecuación~\ref{eq:marco_cos}). El \emph{embedding} de cada clase promedia los \emph{embeddings} del texto de cuatro plantillas, como
  «A dermatoscopic image of \emph{clase}». Luego, se elige la clase de mayor similitud, sin «desconocido».
  \item \textbf{\emph{Linear probe}.} Una regresión logística sobre el \emph{embedding} normalizado, con $C = 1.0$ y pesos
  balanceados por clase.
  \item \textbf{Vecinos más cercanos.} Es el \gls{knn} con $k = 3$ y similitud coseno (Ecuaciones~\ref{eq:marco_knn} y
  \ref{eq:marco_vecinos}). La clase se elige por voto mayoritario de los tres vecinos; si hay empate, gana el vecino
  más similar.
  \item \textbf{SVM con núcleo RBF.} Una \gls{svm} con $C = 10$ sobre \emph{embeddings} estandarizados, con pesos balanceados
  por clase.
\end{enumerate}
Los valores de las cabezas son fijos, es decir, no hubo búsqueda de $C$ ni de $k$ sobre validación.
Sin embargo, con cuatro cabezas sobre la misma representación, un resultado mezcla el efecto de la representación y el
de la cabeza.

Para separar ambos efectos, cabezas clásicas equivalentes (SVM con núcleo RBF y $C = 10$, regresión logística y un
\emph{random forest} de 300 árboles, todas con pesos balanceados) se ajustan sobre el \emph{embedding} de Swin congelado, preentrenado en
\emph{ImageNet}~\cite{liu2021swin}, sin ajuste sobre \emph{HAM10000}. Así, si la diferencia entre las dos
representaciones es la misma con cada cabeza, el efecto es de la representación.

En cuanto a los datos, las cabezas y el control se ajustan sobre el entrenamiento de la partición por lesión
(Sección~\ref{sec:met_particiones}). El modo zero-shot de DermLIP, en cambio, no depende de la partición, porque
no se ajusta con \emph{HAM10000}, y lo mismo vale para los \glspl{vlm} generativos, cuyo ajuste usa \emph{SkinCaRe}
(Sección~\ref{sec:met_generacion}). Todos se evalúan, entonces, sobre la misma prueba que los clasificadores
supervisados. Por último, los \glspl{vlm} generativos, además de clasificar, pueden redactar, y esa es la tarea de la
sección siguiente.
